\begin{figure}[!htbp]
\begin{examplepanel}{srccol}{Source lambda rule}
\small
\[
\src{\inferrule*[right=\textsc{All-I}]{
 d_b:A^\uparrow\cup\{u\};\Gamma,u:T\vdash t:T'\\
 T\ \mathrm{wf}
}{\varnothing;\Gamma\vdash\lambda(u:T).t:((u:T)\to T')^A}}
\]
Only the parameter type and body occur in the source lambda. Its typing
derivation selects the capture annotation and records the body's permitted uses.
\end{examplepanel}
\begin{examplepanel}{tgtcol}{Target lambda with closing evidence}
\[
\tgt{
 \lambda^{\tr{A}}(u:\tr{T}).\
       \underbrace{\tr{d_b}}_{\text{body}}\
       \left\langle
       \underbrace{\tr{d_b}_{\mathrm{uses}}}_{\text{closing evidence }g}
       \right\rangle .}
\]
\small
\[
\tgt{
\tr{\Gamma},u:\tr{T}\vdash^c
 g:\uses{\tr{d_b}}\le_c\tr{A}^{\uparrow}\cup\{u\}.}
\]
\end{examplepanel}
\caption{Translation of a lambda premise in a well-formed source context
\(\Gamma\). The two components are the translated body \(\tr{d_b}\)
and its closing evidence \(\tr{d_b}_{\mathrm{uses}}\).
The superscript \(\uparrow\) carries outer names under the new parameter
binder. Angle brackets delimit stored evidence; the lambda receives only
its term argument at runtime. The functions \texttt{translate} and
\texttt{translateUses} construct these components.}
\label{fig:capture-translation}
\end{figure}
